\section{Finite-rank translated intersections}\label{sec:finite-rank}

We derive Corollary~\ref{cor:finite-rank-translates} from
Theorem~\ref{thm:nondensity} by Pink's reduction
\cite[Theorem~5.3]{Pink2005}, which extends an argument of
R\'emond and Viada \cite[Proposition~4.2]{RemondViada2003} from powers
of an elliptic curve to semiabelian varieties.  The auxiliary ambient
group constructed below is again an abelian variety over $\Qbar$.

For an abelian variety $V$ and an integer $r\ge0$, the set
$V^{[r]}(\Qbar)$ is also the union of $G(\Qbar)$ over all algebraic
subgroups $G\subseteq V$ of codimension at least $r$. Indeed, every
irreducible component of an algebraic subgroup is a torsion coset of
the same dimension. Conversely, a torsion coset $D+\tau$ is contained
in the algebraic subgroup $D+\langle\tau\rangle$, which has dimension
$\dim D$ because $\tau$ is torsion. Thus this notation agrees with the
algebraic-subgroup union used in Pink's reduction.

\begin{proof}[Proof of Corollary~\ref{cor:finite-rank-translates}]
Put $r=\dim_{\mathbb Q}(\Gamma\otimes_{\mathbb Z}\mathbb Q)$ and choose
$a_1,\ldots,a_r\in\Gamma$ whose images form a basis of this vector
space. Let $D\subseteq A^r$ be the Zariski closure of the subgroup
generated by $(a_1,\ldots,a_r)$. Choose a positive integer $m$ that
annihilates the finite component group $D/D^0$. Multiplication by $m$
is surjective on the abelian variety $D^0$, so $[m]D=D^0$. It is also
a finite, hence closed, morphism on $A^r$. Consequently
\[
 \overline{\langle (ma_1,\ldots,ma_r)\rangle}=[m]D=D^0.
\]
Replace each $a_i$ by $ma_i$, write $a=(a_1,\ldots,a_r)$, and put
$C=D^0$. The images of the new $a_i$ still form a basis of
$\Gamma\otimes_{\mathbb Z}\mathbb Q$, and the cyclic subgroup generated
by $a$ is Zariski dense in the connected abelian subvariety $C$.
When $r=0$, this construction uses the empty tuple and the zero abelian
variety. For every $\gamma\in\Gamma$ there are integers $n>0$ and
$n_1,\ldots,n_r$ such that
\begin{equation}\label{eq:finite-rank-relation}
 n\gamma=n_1a_1+\cdots+n_ra_r.
\end{equation}
Indeed, a relation after tensoring with $\mathbb Q$ has only a torsion
error, which a further positive multiple kills.

Set $B=A\times C$ and $Y=X\times\{a\}$. Both are defined over $\Qbar$;
$B$ is abelian and $Y$ is irreducible of dimension $d$. We check that
$Y$ is contained in no proper algebraic subgroup of $B$. If an
algebraic subgroup $J\subseteq B$ contains $Y$, fix $x_0\in X(\Qbar)$.
Then $J$ contains $(x-x_0,0)$ for every $x\in X(\Qbar)$. The Zariski
closure of the subgroup of $A$ generated by these differences is all of
$A$: otherwise its translate by $x_0$ would be a proper translate
containing $X$.
It follows that $A\times\{0\}\subseteq J$. The projection of $J$ to
$C$ is a closed algebraic subgroup containing $a$, and hence is $C$.
These two facts imply $J=B$. Since every proper torsion coset lies in a
proper algebraic subgroup of the same dimension, $Y$ is contained in
no proper torsion coset. Theorem~\ref{thm:nondensity} therefore shows
that $Y(\Qbar)\cap B^{[d+1]}(\Qbar)$ is not Zariski dense in $Y$.

It remains to verify the inclusion
\begin{equation}\label{eq:finite-rank-inclusion}
 \bigl(X(\Qbar)\cap(A^{[d+1]}(\Qbar)+\Gamma)\bigr)\times\{a\}
 \subseteq Y(\Qbar)\cap B^{[d+1]}(\Qbar).
\end{equation}
Let $x=g+\gamma$ belong to the intersection on the left, with
$g\in G(\Qbar)$ for an algebraic subgroup $G\subseteq A$ of codimension
at least $d+1$ and $\gamma\in\Gamma$. Choose $n,n_i$ as in
\eqref{eq:finite-rank-relation}, and define homomorphisms
\[
 \varphi:A^r\longrightarrow A,\qquad
 (z_1,\ldots,z_r)\longmapsto\sum_i n_i z_i,
 \qquad
 \psi:C\longrightarrow B,\qquad
 c\longmapsto(\varphi(c),nc).
\]
The algebraic subgroup
\[
 H=[n]_B^{-1}\bigl((G\times\{0\})+\psi(C)\bigr)
\]
contains $(x,a)$, since
\[
 [n]_B(x,a)=(ng+n\gamma,na)=(ng,0)+\psi(a).
\]
Here $[n]_B^{-1}$ denotes preimage under multiplication by $n$ on $B$.
The kernel of $\psi$ is contained in $C[n]$, and so is the preimage in
$C$ of $\psi(C)\cap(G\times\{0\})$. Thus both the kernel and this
intersection are finite. Multiplication by $n$ on $B$ is finite as well.
Therefore
\[
 \dim H=\dim G+\dim C,
 \qquad
 \codim_B H=\codim_A G\ge d+1.
\]
The algebraic-subgroup description of $B^{[d+1]}$ now gives
$(x,a)\in B^{[d+1]}(\Qbar)$, proving
\eqref{eq:finite-rank-inclusion}. The isomorphism $X\longrightarrow Y$,
$x\longmapsto(x,a)$, and the non-density already proved for the right
side yield the asserted non-density in $X$. When $d=1$, the intersection
is contained in a proper closed subset of the irreducible curve $X$,
and every such subset is finite.
\end{proof}

The algebraicity of the translating points is used to place $Y$ over
$\Qbar$, where Theorem~\ref{thm:nondensity} applies. Thus this deduction
is the abelian algebraic-point case of Pink's Conjecture~5.2, not its
general formulation for semiabelian varieties over $\mathbb C$. In
higher dimensions the conclusion is non-density, and the proper closed
exceptional subset may have positive dimension.
